\chapter{Beta-D-Glucan}

\section*{What Is This Test?}
Beta-D-glucan is a cell-wall component released by many fungi.
Serum beta-D-glucan testing is a supportive biomarker.
It is used for selected invasive fungal infections.
These include invasive candidiasis, aspergillosis, and Pneumocystis jirovecii pneumonia.

\section*{Alternative Names}
\begin{itemize}
  \item BDG
  \item 1,3-Beta-D-Glucan
  \item Fungal Biomarker Test
\end{itemize}
\section*{Why This Test Is Ordered}
Testing is used in selected high-risk patients.
It is used when an invasive fungal infection is possible.
Such patients include those with prolonged neutropenia, hematologic malignancy, transplant, intensive-care illness, or unexplained pneumonia.
The test is intended to complement culture, antigen, molecular, and clinical assessment.
It does not replace them.

\section*{How This Test Is Performed}
A blood sample is collected into a specialized tube.
The laboratory analyzes the sample with a sensitive colorimetric or kinetic assay.
Serial measurements may be more informative than a single result in high-risk patients.

\section*{Preparation, Risks, and Limitations}
No fasting is required.
Venipuncture is the usual risk.
Several factors can cause false-positive results.
These include hemodialysis, albumin or immunoglobulin products, some antibiotics, gauze exposure, severe mucositis, and sample contamination.
The test does not reliably detect Mucorales or Cryptococcus.
It does not identify the organism.

\section*{Understanding Results}
A positive result increases suspicion.
It is not diagnostic by itself.
A negative result lowers the likelihood of some invasive fungal infections.
It cannot exclude disease.
Antifungal treatment is one such example.
Poorly detected organisms are another.

\section*{References}
European Organisation for Research and Treatment of Cancer and Mycoses Study Group guidance on invasive fungal disease definitions.

\clearpage
\section*{Understanding Results and Follow-up}
The laboratory report should identify the specimen, method, units, reference interval or decision limit, and whether the result is positive, negative, abnormal, or inconclusive. Reference intervals vary with age, sex, pregnancy, specimen type, assay, and laboratory; a result outside a range is not automatically a diagnosis, and a result within a range does not always exclude disease. Discuss unexpected results with the ordering clinician, who can decide whether repeat or confirmatory testing, treatment, or referral is appropriate. Do not change medicines or delay urgent care based on an isolated result.


\section*{Regulatory Status and Authorization Date}
This chapter describes a test category, not a specific commercial product. FDA, CDSCO, EU IVDR, and MHRA approval or registration dates therefore cannot be assigned without the assay manufacturer, product name, instrument, and intended use. For laboratory-developed tests, an individual agency approval date may not exist. See the Preface for the interpretation of regulatory status.

\clearpage
